\section{Introduction}

For each $m\in\Z_{>0}$, consider the integer recurrence
\begin{equation}
 b_1=m,\qquad b_{n+1}=b_n+(b_n\bmod n).
 \label{eq:recurrence}
\end{equation}
The recurrence appears in OEIS A073117 and A117846 and was discussed in MathOverflow question 191518; see \cite{oeisA073117,oeisA117846,mathoverflow191518}. Cloitre's stabilization conjecture asks whether every positive starting value $m$ produces an orbit whose increments $b_{n+1}-b_n$ are eventually constant.

The conjecture itself is not resolved here. Our subject is the smaller and logically independent question of which constants can occur when stabilization does happen.

\begin{definition}[Attainable eventual-increment spectrum]
Let $\Spec$ be the set of positive integers $c$ for which there exist $m\ge1$ and $t\ge1$ such that
\[
 b_{n+1}-b_n=c\qquad\text{for all }n\ge t.
\]
\end{definition}

Abercrombie asked in OEIS A117846 whether these values include every positive integer \cite{oeisA117846}. Our main result answers this negatively.

\begin{theorem}[Main theorem]
The attainable spectrum $\Spec$ is not equal to $\Z_{>0}$. More precisely,
\[
 1,2,3,4,6\in\Spec,
 \qquad
 5,7\notin\Spec.
\]
Consequently, $5$ and $7$ are the two smallest omitted eventual increments.
\end{theorem}

The proof has two distinct components. The mathematical component is an effective bound: if a start $m$ stabilizes with increment $c$, then $m<(c+3)(3c+5)$. The computational component is a complete exact certificate for all $1\le m<260$. The certificate is sufficient because the theorem reduces the cases $c=5$ and $c=7$ to $m<160$ and $m<260$, respectively. No assumption is made about starts outside these finite intervals or about universal stabilization.
